\chapter{Dental Crowns}

\section*{Overview}
A dental crown is a custom-made cap placed over a damaged or weakened tooth to restore its shape, strength, function, and appearance. The exact approach, setting, and preparation depend on the indication, anatomy, and the patient's health.

\section*{Why It Is Done}
Crowns may be used for extensive decay, a fractured tooth, large fillings, root-canal-treated teeth, worn teeth, or to restore an implant.

\section*{How It Is Performed}
The tooth is examined and prepared, an impression or digital scan is taken, and a temporary crown may be placed. The permanent crown is later cemented or bonded after fit and bite are checked.

\section*{Risks and Follow-up}
Sensitivity, bite problems, loosening, fracture, recurrent decay, gum irritation, and need for replacement are possible. Brush, floss, and attend routine dental visits.

\section*{Outcome and Follow-up}
The expected outcome depends on the reason for the procedure, the person's health, and whether additional treatment is needed. The clinician reviews the result with the patient, explains any next steps, and sets follow-up according to the findings and recovery.

\section*{When to Seek Medical Care}
Follow the procedure team's specific instructions. Seek urgent medical assessment for severe or worsening pain, uncontrolled bleeding, fever or signs of infection, trouble breathing, chest pain, fainting, new weakness or confusion, inability to urinate, or any rapidly worsening symptom. The appropriate warning signs depend on the procedure and the patient's medical condition.

\section*{References}
\begin{enumerate}
  \item MedlinePlus. Dental Crowns. U.S. National Library of Medicine; 2025.
  \item Current Medical Diagnosis and Treatment. McGraw-Hill; 2025.
\end{enumerate}
\clearpage
